% sections/05-discussion.tex

\section{Discussion}\label{sec:discussion}

The energy functionals introduced in Section~\ref{sec:methods} are
descriptive tools: they characterize the reconfiguration landscape but
do not, by themselves, prove that safe paths exist. Their value lies
in revealing \emph{structure} — connections to established mathematical
frameworks that may yield rigorous results. We discuss four such
connections.

\subsection{Connection to TQFT and the Penrose Evaluation}
\label{sec:tqft}

The tetrahedral embedding (Section~\ref{sec:tet-embed}) maps the four
colors to vectors in $\mathbb{R}^3$ that span the spin-1
(adjoint) representation of $\mathfrak{so}(3)$. In this representation:

\begin{itemize}
  \item A proper 4-coloring of a planar triangulation corresponds to a
    spin network state — an assignment of $\mathfrak{so}(3)$
    representations to the edges of the dual graph.
  \item A Kempe $(a,b)$-swap corresponds to a $6j$-symbol recoupling:
    the $180^\circ$ rotation (Proposition~\ref{prop:kempe-rotation})
    exchanges two basis vectors, which is analogous to the action of a
    recoupling coefficient in the Turaev-Viro state
    sum~\cite{turaev1992} (though the correspondence has not been made
    rigorous).
  \item The Penrose evaluation $\Pen(G)$~\cite{kauffman1990} of a
    bridgeless planar trivalent graph counts the number of proper
    3-edge-colorings (equivalently, proper 4-face-colorings of the
    dual), weighted by signs. Kauffman proved that $\Pen(G) \neq 0$
    is equivalent to the 4CT.
\end{itemize}

The Magic Gem covariance energy provides a \emph{continuous relaxation}
of the Penrose state sum. While $\Pen(G)$ counts discrete colorings
with signs, $E_{\mathrm{MG}}(G, c)$ measures how far a given coloring
is from the balanced (zero-centroid) condition that characterizes
contributions to $\Pen(G)$. It is tempting to conjecture that the
two-basin structure observed in Section~\ref{sec:kinetic-trap}
corresponds to the sign structure in the Penrose evaluation —
Basin~A to opposite-sign contributions, Basin~B to same-sign — but we
have no direct evidence for this correspondence. If it can be made
precise, it would connect the kinetic trap structure to the
nonvanishing of $\Pen(G)$.

\begin{remark}\label{rem:tqft-rigidity}
We speculate that surface tension rigidity
(Definition~\ref{def:st-rigidity}) might characterize degeneracy in the
TQFT state space: zero-variance surface tension would correspond to
Kempe chains that are rigid under the $\mathrm{SO}(3)$ action. However,
the Surface Tension Rigidity Conjecture was falsified at $n \geq 10$
(Remark~\ref{rem:st-falsified}), so any such characterization would
need to be substantially revised.
\end{remark}

\subsection{Connection to Sheaf Cohomology}\label{sec:sheaf}

Cellular sheaves on graphs~\cite{ghrist2014} provide a cohomological
framework for constraint satisfaction. We outline a potential
formalization here; no results in this paper depend on it, and we do not
establish any sheaf-theoretic theorems. Given a graph $G = (V, E)$ and
a proper $k$-coloring $c$, define a sheaf $\mathcal{F}_c$ as follows:
\begin{itemize}
  \item To each vertex $v$, assign the stalk $\mathcal{F}_c(v)$: the
    set of colors available to $v$ given the coloring of its neighbors.
  \item To each edge $(u, v)$, assign the restriction map: the
    constraint that $c(u) \neq c(v)$.
\end{itemize}
The zeroth cohomology $H^0(G, \mathcal{F}_c)$ encodes global colorings
consistent with the local constraints. The first cohomology
$H^1(G, \mathcal{F}_c)$ measures obstructions to extending local
colorings globally.

We note one suggestive (but unproven) analogy with surface tension:
\begin{itemize}
  \item The boundary edges of a Kempe chain $K$ are precisely the
    edges in the \v{C}ech nerve of the sheaf cover restricted to the
    $(a,b)$-subgraph.
  \item Rigidity ($\Gamma_{ab} = 0$) means all chains have the same
    normalized boundary. We speculate that this corresponds to maximally
    constrained restriction maps and hence nonzero $H^1$, but we have
    not verified this correspondence, and the Surface Tension Rigidity
    Conjecture itself was falsified at $n \geq 10$.
  \item Motivated by the $n = 9$ data (where the correspondence held),
    we tentatively define: a coloring sheaf has \emph{planar-positivity}
    if no Kempe chain in any color pair has rigid surface tension.
    Whether planar-positivity implies trivial $H^1$ obstructions to
    reconfiguration — let alone safe path existence — remains entirely
    open.
\end{itemize}

\subsection{Connection to Chromatic Polynomials}
\label{sec:chromatic}

The Potts partition function
(Section~\ref{sec:potts}) provides a bridge between the energy
landscape and the chromatic polynomial:
\begin{equation}\label{eq:potts-bridge}
  P(G, q) = \lim_{\beta \to \infty} Z(G, q, \beta).
\end{equation}
The chromatic polynomial evaluates via deletion-contraction:
$P(G, q) = P(G - e, q) - P(G / e, q)$ for any edge $e$. Each term in
this recursion corresponds to a ``channel'' in the reconfiguration
graph — either the edge $e$ is bichromatic (deletion) or monochromatic
(contraction).

We speculate that the two-basin structure in the energy landscape
(Section~\ref{sec:kinetic-trap}) may correspond to the sign structure
in deletion-contraction — Basin~A to contraction channels where the
edge constraint is tight, Basin~B to deletion channels where the
constraint is relaxed. This analogy is suggestive but unverified: we
have not shown that individual basins align with specific
deletion-contraction terms.

\begin{remark}
The chromatic polynomial $P(G, q)$ has roots on the real line that
accumulate at $q = 4$ for planar graphs (the Beraha conjecture). Whether
this accumulation point is related to the kinetic trap structure observed
in Section~\ref{sec:kinetic-trap} is an open question; we have no
evidence connecting the two.
\end{remark}

\subsection{Connection to Discharging}\label{sec:discharging}

The discharging method — used in both proofs of the
4CT~\cite{appel1977, robertson1997} — assigns charges to vertices and
faces, then redistributes charge via local rules to derive
contradictions. The electrostatic potential
(Section~\ref{sec:electrostatic}) generalizes this:

\begin{itemize}
  \item The initial charge assignment $\rho(v) = \delta_{c(v), 5}$ is
    analogous to the standard initial charge $6 - \deg(v)$ in
    discharging, but defined on \emph{colorings} rather than on the
    graph structure.
  \item The Laplacian diffusion $V = L^+ \rho$ is a continuous analog
    of the discrete charge redistribution rules.
  \item The electric field gradient $\nabla E(v)$ identifies vertices
    where the charge distribution is most uneven — potential locations
    for discharging rules that resolve the fifth color.
\end{itemize}

More speculatively, surface tension was initially hypothesized to guide
the search for smaller unavoidable sets. However, the falsification of
the Surface Tension Rigidity Conjecture at $n \geq 10$
(Remark~\ref{rem:st-falsified}) undermines this idea: rigid surface
tension does not reliably identify the hard cases at larger scales. A
revised version of this speculation would need a different indicator —
perhaps the local entropy or kinetic trap signatures that do persist.

Separately, the Magic Gem energy might serve as a \emph{reducibility
criterion}: the Magic Gem energy of a boundary coloring, viewed as a
function of possible interior extensions, could in principle detect
non-extendability via high energy barriers. We have not tested this
idea computationally.

\subsection{Limitations}\label{sec:limitations}

We acknowledge several limitations of the present approach:

\begin{enumerate}
  \item \textbf{Descriptive, not prescriptive.} The energy functionals
    characterize the landscape but do not directly produce safe paths.
    A prescriptive version would require proving that gradient descent
    on, say, $E_{\mathrm{MG}}$ with a ridge-crossing heuristic always
    finds a safe path.

  \item \textbf{Continuous vs.\ discrete.} The tetrahedral embedding
    and Laplacian potential are continuous objects defined on a discrete
    graph. The continuous approximation cannot capture all discrete
    merge constraints. In particular, the surface tension rigidity
    condition (Definition~\ref{def:st-rigidity}) is a
    \emph{combinatorial} statement that the continuous energy
    functionals can only approximate.

  \item \textbf{Overlapping distributions.} While local entropy and
    surface tension rigidity discriminate well on the $n = 9$
    counterexamples, the distributions of safe and unsafe values
    overlap for all individual metrics. A composite score combining
    multiple functionals would likely improve discrimination, but we
    have not optimized such a combination.

  \item \textbf{Small sample size for energy analysis.} The energy
    functional analysis in
    Sections~\ref{sec:discriminating}--\ref{sec:kinetic-trap} is
    verified only at $n = 9$ (48 counterexample colorings across 2
    triangulations). Subsequent enumeration at $n = 10$ tested
    approximately 1.9 million cases and found safe paths in 100\% of
    instances, with maximum safe detour cost
    $d_{\text{safe}} - d_{\text{opt}} = 1$, providing further evidence
    for safe path existence but on a broader dataset than the energy
    functional analysis was originally applied to.

  \item \textbf{Surface Tension Rigidity Conjecture falsified.}
    The Surface Tension Rigidity Conjecture
    (Conjecture~\ref{conj:st-rigidity}) was falsified by subsequent
    computation at $n \geq 10$. While zero surface tension rigidity
    correlates with merge-proneness at $n = 9$, this correlation breaks
    down at larger scales. The conjecture as stated is false; however,
    the local entropy discrimination and kinetic trap structure continue
    to hold in the extended dataset.
\end{enumerate}
